\section{Actor-critic}


In this problem we consider a realistic task in which a Sawyer robot needs to grasp a tool and use it to hammer in a nail. 
The agent controls a 4-DOF Sawyer robot with a continuous action space and receives environment states as observations.  
Providing dense rewards for real world problems is difficult, so in this environment the agent receives a sparse reward of 1.0 upon fully completing the task and no intermediate rewards. 
We will implement a general actor-critic algorithm that learns to solve the task with almost 100\% success rate. 
At each step the agent processes observations $o_t$, which consists of the environment states (e.g., pose of the robot gripper and hammer, etc.) from the last two environment steps. 
The full agent consists of:

\begin{itemize}
  \item The actor-critic algorithm utilizes critic networks $Q_{\theta^i}(o_t, a_t)$. This is implemented as the \textbf{Critic} class in \textbf{submission.py}. We use an ensemble 
  of $N$ different critic networks. Using an ensemble of critic networks can help reduce error in the Q-value estimates.
  \item Each critic network has a corresponding $target$ critic network $\bar{Q}_{\theta^i}(o_t, a_t)$., which is used for computing the Bellman targets for critic updates. 
  We use target critic networks, which are updated more slowly, to have a more stable target for learning.
  \item The final component of the algorithm is the Actor or policy $\pi_{\theta}(o_{t})$. This is implemented as the \textbf{Actor} class in \textbf{submission.py}. 
  You should familiarize yourself with these functions and their inputs and outputs. You should not modify any other files besides \textbf{submission.py}.
\end{itemize}


\begin{enumerate}[(a)]

  \item \points{2a} 

Optimizing behavior in environments with sparse rewards is difficult due to limited reward supervision. To alleviate this, we provide the agent with 20 successful demonstrations, which we will use to pre-train with behavior cloning. In \textbf{submission.py} complete the \textbf{bc} function of the \textbf{ACAgent} class to train the policy using supervised behavior cloning. That is, given state-action pairs $o_t, a_t$ optimize the loss $$\mathcal{L}_{\pi_{\theta}, f_{\theta}}(o_t, a_t) = -\log \pi_{\theta}(a_t|o_t)$$

We will also use this to update the actor during later RL training, to improve stability.


  \input{02-actor-critic/02-actor-critic}

  \input{02-actor-critic/03-run}

\end{enumerate}



